\documentclass[a4paper,11pt]{article}
\usepackage[utf8]{inputenc}
\usepackage[T1]{fontenc}
\usepackage[french]{babel}
\usepackage[left=2cm,right=2cm,top=2cm,bottom=2cm]{geometry}
\usepackage{xcolor}
\usepackage{hyperref}
\usepackage{fontawesome5}
\usepackage{parskip}

% --- Couleurs Eiffage Rail ---
\definecolor{primary}{RGB}{30, 90, 50}
\hypersetup{colorlinks=true, urlcolor=primary}

\begin{document}

% --- EXPÉDITEUR ---
\noindent
\begin{minipage}[t]{0.5\textwidth}
    \textbf{Loïc Aron MBASSI EWOLO}\\
    Élève-Ingénieur Bac+5 -- Double diplôme ENSEA\\[3pt]
    \faMapMarker\ Cergy, France (Mobile IDF)\\
    \faPhone\ (+33) 7 59 52 33 98\\
    \faEnvelope\ \href{mailto:loicaron.mbassiewolo@ensea.fr}{loicaron.mbassiewolo@ensea.fr}
\end{minipage}
\hfill
% --- DATE ---
\begin{minipage}[t]{0.4\textwidth}
    \raggedleft
    Cergy, le 27 février 2026
\end{minipage}

\vspace{1.2cm}

% --- DESTINATAIRE ---
\hfill
\begin{minipage}[t]{0.5\textwidth}
    \raggedleft
    \textbf{Eiffage Rail}\\
    Service Recrutement\\
    Élancourt, France
\end{minipage}

\vspace{1cm}

\noindent\textbf{Objet :} Candidature -- Alternance / Stage PFE Informaticien F/H (Développement Logiciel Applicatif)

\vspace{0.8cm}

Madame, Monsieur,

Maintenir et faire évoluer des applications métier utilisées au quotidien par des équipes terrain, c'est une mission que je connais concrètement : pendant deux ans, j'ai assuré le développement et la maintenance d'une plateforme de télémédecine complète pour SOSDOCTA (Belgique/Canada), incluant des dashboards de suivi en temps réel, des systèmes d'alertes, et la gestion de données sensibles. Cette expérience de durée et de responsabilité -- rare à ce stade de formation -- me semble directement transférable aux applications de gestion de matériel et de chantier décrites dans votre offre.

La logique de traçabilité et d'alertes que j'ai implémentée chez SOSDOCTA (suivi des rendez-vous, historique patients, alertes de statut) est structurellement identique à celle de votre outil de gestion de matériel : visites périodiques, agréments, entretiens et alertes associées. Cette correspondance n'est pas anecdotique -- elle implique les mêmes choix de conception de base de données, de gestion des états et de notifications. Sur le plan technique, j'ai consolidé cette expérience avec un poste de développeur backend en fintech (CSC, Laravel, MySQL, Docker, Redis) et la conception d'une architecture microservices utilisant MongoDB (NoSQL) pour la gestion de données non structurées -- deux contextes qui couvrent les compétences SQL/NoSQL mentionnées dans l'offre.

Ma pratique en Java, illustrée par un projet applicatif complet utilisant les patterns Singleton, Factory et Observer, renforce ma capacité à intervenir rapidement sur des bases de code existantes et à contribuer à leur évolution sans introduire de dette technique. L'habitude du travail en équipe -- coordination de cellules projets à l'ENSPY, reporting régulier lors de missions freelance, collaboration à distance sur plusieurs fuseaux horaires -- me permet de m'intégrer efficacement dans une organisation décentralisée comme la DOP Rail.

Je suis disponible dès avril 2026 pour un stage PFE, ou dès septembre 2026 pour une alternance, sur le site d'Élancourt. La perspective de contribuer à des outils qui s'ancrent directement dans l'organisation des chantiers ferroviaires, avec des effets mesurables sur l'efficacité opérationnelle, représente exactement le type de mission dans laquelle je souhaite m'investir.

Je reste à votre disposition pour un entretien.

Je vous prie d'agréer, Madame, Monsieur, l'expression de mes salutations distinguées.

\vspace{0.8cm}

\textbf{Loïc Aron MBASSI EWOLO}\\
\textit{Élève-Ingénieur ENSEA / Polytechnique Yaoundé}

\end{document}
